\documentclass[10pt,a4paper]{article}
\usepackage[margin=25mm]{geometry}
\usepackage{amsmath,amssymb,lmodern}
\setlength{\emergencystretch}{2em}
\title{Orthogonal complements of adjoint invariant subspaces\\\large Proof of Conjecture 00000006837}
\author{ziangni-sys}
\date{4 October 2026}
\begin{document}
\begin{center}
{\Large Orthogonal complements of adjoint invariant subspaces}\\[4pt]
{\large Proof of Conjecture 00000006837}\\[6pt]
ziangni-sys \qquad 4 October 2026
\end{center}
\section*{Precise statement and convention}
The source pairs invariant subspaces of an operator and its adjoint
through orthogonal complements. We prove this standard correspondence
for every bounded linear operator $A$ on any real or complex Hilbert
space $H$, without a finite-dimensional assumption.
Write $A^*$ for the Hilbert-space adjoint, characterized by
$\langle x,A^*y\rangle=\langle Ax,y\rangle$, and define
\[
U^\perp=\{y\in H:\langle x,y\rangle=0\text{ for every }x\in U\}.
\]
An invariant subspace means $A(U)\subseteq U$.
In the invariant-subspace lattice of a Hilbert-space operator,
subspaces are conventionally closed. We also prove the forward
implication for arbitrary linear subspaces, and identify explicitly
where closedness is needed for the converse.

\section*{Adjoint invariance}
Suppose $A(U)\subseteq U$. For $y\in U^\perp$ and $x\in U$,
\[
\langle x,A^*y\rangle=\langle Ax,y\rangle=0,
\]
because $Ax\in U$. Thus $A^*(U^\perp)\subseteq U^\perp$.
No closedness of $U$ was used in this direction.
Conversely, suppose $U$ is closed and $A^*(U^\perp)\subseteq U^\perp$.
Apply the result just proved to $A^*$ and $U^\perp$.
Using $(A^*)^*=A$ and $(U^\perp)^\perp=U$ gives $A(U)\subseteq U$.
Therefore, for every closed $U$,
\[
 A(U)\subseteq U \quad\Longleftrightarrow\quad
 A^*(U^\perp)\subseteq U^\perp.
\]
For a nonclosed subspace, the double orthogonal complement is
$\overline U$, not necessarily $U$, so this converse is not asserted
without closedness.

\section*{Full correspondence and order}
Every orthogonal complement is closed. Consequently
$U\mapsto U^\perp$ maps closed $A$-invariant subspaces to closed
$A^*$-invariant subspaces. Its inverse is the same orthogonal-complement
operation in the other direction, since both double-adjoint and
double-complement identities hold.
For any closed subspaces $U,V$,
\[
U\subseteq V\quad\Longleftrightarrow\quad V^\perp\subseteq U^\perp.
\]
The forward implication follows from the definition of orthogonality;
the reverse follows by taking orthogonal complements again.
Thus this is an order-reversing bijection between the two closed
invariant-subspace lattices. This establishes the orthogonal-complement
pairing claimed in both source versions, in the standard bounded
Hilbert-space adjoint setting.

\section*{Formal verification}
Lean is generic over a scalar field with \texttt{RCLike} structure
and a complete inner-product space, hence covers both $\mathbb R$
and $\mathbb C$. It uses actual Mathlib continuous linear operators,
their adjoints, submodules, orthogonal complements, and topological
closedness. The predicate of invariance is the actual condition
$\forall x\in U,\ Ax\in U$.
The project proves the arbitrary-subspace implication, the closed
subspace equivalence, a bundled equivalence of the two types of closed
invariant subspaces, and the inclusion-reversal equivalence.
Its axiom audit uses only standard logical axioms; no analytic
correspondence or invariance hypothesis is silently assumed as a result.
\end{document}
